\documentclass[11pt,a4paper]{amsart} % Or article/book

% --- GLOBAL PREAMBLE ---
\usepackage{amsmath,amssymb,amsfonts,amsthm}
\usepackage{mathtools}
\usepackage{graphicx}
\usepackage{hyperref}
\usepackage{enumitem}

% --- THEOREM STYLES & COMMANDS ---
\newtheorem{theorem}{Theorem}[section]
\newtheorem{exercise}{Exercise}[section]
\newenvironment{solution}{\begin{proof}[Solution]}{\end{proof}}

\newcommand{\R}{\mathbb{R}}
\newcommand{\C}{\mathbb{C}}

\begin{document}

    \title{Textbook Solutions Manual}
    \author{Your Name}
    %	\maketitle

    % --- INTERACTIVE COMPILATION ZONE ---
    % Comment out lines using '%' to only preview what you are currently testing

    %\input{./Appendix/Product of subgroups of a group.tex}
    \section{Exercise 2}
    Let $A$ be a ring and let $A[x]$ be the ring of polynomials in an indeterminate $x$, with coefficients in $A$. Let $f=a_0+a_1x+\cdots+a_nx^n\in A[x]$.

    \textbf{Question 2.i}

    Prove that $f$ is a unit in $A[x]$ $\Leftrightarrow$ $a_0$ is an unit in $A$ and $a_1,\ldots,a_n$ are nilpotent.

    \textbf{Hint:} If $b_0+\cdots b_mx^m$ is the inverse of $f$, prove by induction on $r$, that $a_n^{r+1}b_{m-r}=0$. Hence show that $a_n$ is nilpotent and then use Ex. $1$.

    \textbf{Solution} Suppose we use the hint, then we have to prove that dfd

\end{document}